\section*{章节综合：一次看懂反复作用}
把本章的概念放进同一个问题：矩阵 $A=\left[\begin{smallmatrix}2&1\\1&2\end{smallmatrix}\right]$ 反复作用于 $(3,1)$，怎样求第 $n$ 次结果？先不要连乘矩阵。寻找被保留的方向 $u=(1,1)$、$v=(1,-1)$，它们的倍率分别为三和一，而且 $(3,1)=2u+v$。

因此 $A^n(3,1)=2\cdot3^n u+v=(2\cdot3^n+1,2\cdot3^n-1)$。取 $n=2$ 得 $(19,17)$，可用两次普通矩阵乘法交叉检查。换基把复合变成了对每份坐标单独乘倍率，特征向量的用途就在这里。

再问这个动作会丢失信息吗？两特征倍率都非零，行列式为三，面积扩大三倍且保持定向，所以可以求逆。若把第二倍率改为零，则所有沿 $v$ 的变化都不可见，原来唯一的坐标恢复变成多解问题。这样，基、矩阵复合、行列式、核与特征值就在一个例子中汇合了。

完成本章后，试着分别用一幅图、一段文字和一组坐标式解释这个问题。三种表达应讲述同一个动作。这里的综合推导供纸笔检查；机器证明仍以各课列出的实际定理为准。
